\section{法的本体}
\begin{enumerate}
  \item 做题，\zd{抠字眼讲逻辑}。关键是听懂，理解，记下来是最没用的。
  \item 法概念争议。恶法亦法 vs 恶法非法，恶指道德上的邪恶。
        \begin{enumerate}
          \item 争议前提。法律和道德是否\zd{存在概念（本质）上的必然联系}。
          \item 实证主义。
                \begin{enumerate}
                  \item \zd{权威性制定} \yjt 分析法学。
                  \item \zd{社会实效}（认同程度）\yjt 法社会学/法律现实主义。
                \end{enumerate}
          \item 非实证主义。一回事。
                \begin{enumerate}
                  \item \zd{内容的正确性 = 道德正确} \yjt 古典自然法。
                  \item 3 要素全都看 \yjt 综合法学派。
                \end{enumerate}
        \end{enumerate}
  \item 马列。官方正统。3 性，有层次差异。3 位一体，构成区别。
        \begin{enumerate}
          \item 阶级分析法。经济基础，上层建筑。
          \item 正式/国家/官方性。
          \item 阶级性。体现统治阶级的意志。代表最先进的生产力。意志 3 性。
                \begin{enumerate}
                  \item 唯一性。是、并且永远是统治阶级意志的体现。
                  \item 整体性、长远性。整体、长远利益。
                \end{enumerate}
          \item \zd{社会性}。物质制约性。物质生活条件（硬件）才是法律（软件）内容的\zd{决定性}因素。最终非唯一。
        \end{enumerate}
  \item 国法 = 国家现行法。法概念的\zd{交集}。大家都认可的部分。
        \begin{enumerate}
          \item 制定法/成文法。
          \item 判例法、习惯法 \yjt 不成文法。
          \item \ydy 法律，发挥法律职能。党纪，教会。
        \end{enumerate}
  \item 特征。对个性的描述。为了区分。
        \begin{enumerate}
          \item 调整\zd{人行为}的社会（涉他）规范。\zd{不针对思想、自我行动}。法律是最底限的道德。
          \item 国家意志性。
                \begin{enumerate}
                  \item 从无到有的创造。内容
                  \item 认可（明示、默示）从无效到有效的转化。效力
                \end{enumerate}
          \item 普遍性。\zd{国家权力所及}的范围内。
          \item 以\zd{权利和义务}为内容的社会规范。一致性。
          \item 国家强制。
          \item 程序性。
          \item 可诉性。宪法除外。
        \end{enumerate}
  \item 作用 = 影响。分类。
        \begin{enumerate}
          \item 对个体。规范作用。基础和前提。必考。
                \begin{enumerate}
                  \item 指引。\zd{对象是否特定}。未发生。\al{依法办事}。
                        \begin{enumerate}
                          \item 个别性/非规范性。
                          \item 规范性指引。\yi 确定性看义务。应为、勿为。 \er 不确定性看权利（可放弃）。可为。
                        \end{enumerate}
                  \item 评价。针对他人。他人已发生，合法与违法。\al{判决}。
                  \item 教育。示警（反面）、示范（正面）。\al{举例子}。针对一般人，必然已发生。
                  \item 预测。现实法律关系为前提。
                  \item 强制。违法犯罪已发生。
                \end{enumerate}
          \item 对社会。社会作用。结果和目的。维护阶级统治、社会正常存在。
          \item 局限性。\bdy 缺点、缺陷（可以克服、改正）。无法突破的。\al{是人就飞不起来}。
                \begin{enumerate}
                  \item 滞后、模糊、僵硬、空白、漏洞。法律万能主义。不要强法所难、其他手段不可或缺。
                \end{enumerate}
        \end{enumerate}
  \item 价值。客体对主体需求的满足。主客体相关的概念。\al{沙漠里给你水，给你讲义}。
        \begin{enumerate}
          \item 能不能满足人的需求；满足哪些需求；冲突处理。
          \item 自由。受限的自由。
                \begin{enumerate}
                  \item 伤害原则。禁止伤害他人，侵害他人合法权益。
                  \item 家长主义原则。避免自我伤害。\al{开车、系安全带、带头盔、超速}。
                  \item 冒犯原则。对他人情感的伤害。\zd{公众反感}。
                  \item 道德主义原则。\al{见死不救，见义不勇为}。
                \end{enumerate}
          \item 人权。人之为人，应当享受的权利，受到的尊重。
                \begin{enumerate}
                  \item 不依赖其他条件，法律规定。
                  \item 人权法律化。
                \end{enumerate}
          \item 正义、秩序。
          \item 冲突。\zd{矛盾、取舍}。
                \begin{enumerate}
                  \item 价值位阶，保高舍低。具体问题具体分析。\al{吃饭、看电影}。
                  \item 比例原则。成本和限度。两害相权取其轻。
                \end{enumerate}
        \end{enumerate}
  \item 要素。\zd{法律规范 = 法律规则 + 法律原则}。是上位概念。
        \begin{enumerate}
          \item 法律规则。以一定的逻辑结构形式规定法律权利和法律义务的一种法律规范。逻辑结构学说。
                \begin{enumerate}
                  \item 3 要素说、2 要素说。
                  \item 新 3 要素说。考试通说。\zd{假定条件 + 行为模式 + 法律后果}。
                        \begin{enumerate}
                          \item 假定条件。适用前提。主体条件、事项条件。出现/加上\zd{如果}两字。
                          \item 行为模式。表达权利、义务。可为、应为、勿为。定人 + 定事。
                          \item 法律后果。肯定/正面、否定/负面后果。
                        \end{enumerate}
                  \item \zd{逻辑结构要素的省略问题}。
                        \al{故意杀人罪是指，故意杀人的，处死刑、无期徒刑或者十年以上有期徒刑；情节较轻的，处三年以上十年以下有期徒刑}。
                        \begin{enumerate}
                          \item 反推 \hz 其他法条存在。\al{有纳税的义务}。
                          \item \zd{条文中均可省略，逻辑缺一不可}。保持语言简练，体系逻辑。
                        \end{enumerate}
                  \item vs 条文。内容与形式（外在）。随便的关系。\al{水，气态、液态、固态；液体可以是水、硫酸、石油}。
                        \begin{enumerate}
                          \item 规范条文。规则、原则。
                          \item 非规范条文。如表达概念。
                        \end{enumerate}
                  \item vs 语言。
                        \begin{itemize}
                          \item 依赖语言；通过解释达到传到含义；通过语言表达。
                          \item \zd{规范语句}。命令句、允许句。\al{应当 = 必须}。
                          \item 陈述句也可传达法律规则。把道义助动词加进去。
                        \end{itemize}
                  \item 分类。必考。有且只有一个，\zd{分类标准}。
                        \begin{enumerate}
                          \item 授权性和义务性规则。行为模式。应为叫命令，勿为叫禁止，都是义务。
                          \item 确定、委任、准用。内容完整性；找人，找法帮忙。\al{商业银行组织适用公司法的规定}。
                                \begin{itemize}
                                  \item 找人。直接点名、间接点名。\al{找国务院；找行政法规}。
                                  \item 找法。具体名称的法。参考援引现实存在的法。
                                  \item 空白罪状。
                                \end{itemize}
                          \item 强行性、任意性。当事人是否有选择空间。\al{有约从约}。
                        \end{enumerate}
                \end{enumerate}
          \item 法律原则。是规则的基础。\al{里格斯诉帕默尔}。
                \begin{enumerate}
                  \item 最佳化命令（最好的处理效果）。规则是确定性命令。
                  \item 功能。基础、指导。
                  \item 分类。考的少。
                        \begin{enumerate}
                          \item 公理性、政策性。普遍性强弱。放之四海而皆准。\al{诚实守信}。
                          \item 基本原则、具体原则。覆盖面。多领域。\al{法律面前人人平等；罪刑法定}。
                          \item 实体原则、程序原则。\al{原告就被告}。
                        \end{enumerate}
                  \item vs 法律规则。\zd{3 看，内容 \yjt 法条位置 \yjt 部门法}。
                        \begin{enumerate}
                          \item 宽泛抽象。
                          \item \zd{原则能结合个性}。\al{罪刑相适应。不同杀人案量刑不一}。量体裁衣。
                          \item 适用方法。规则 \yjt 全有、全无，互斥，非此即彼。原则 \yjt 兼顾。
                        \end{enumerate}
                  \item 适用条件。\zd{限制自由裁量}。
                        \begin{enumerate}
                          \item 穷尽规则用原则。
                          \item 原则取代规则。\zd{个案正义 + 更强利用（充分论证）}。\al{泸州遗赠二奶案}。
                        \end{enumerate}
                \end{enumerate}
        \end{enumerate}
  \item 权利、义务。
        \begin{enumerate}
          \item 权利法定、权利自主、共生性。
          \item 义务。将做还未做的行动，不能放弃。\al{还钱}。
          \item 分类。
                \begin{enumerate}
                  \item 基本、普通。宪法规定。
                  \item 绝对、相对。物权，债权。
                  \item 个体、集体、国家。
                \end{enumerate}
        \end{enumerate}
  \item 渊源和分类。水的来源、出处。找法。\zd{效力}。\zd{定性}（有无之分）、\zd{程度}（高下之分）。
        \begin{enumerate}
          \item 法律渊源、非法律渊源。
          \item 正式。明文规定，必须考虑。参考立法法。\zd{法律解释没有独立性}，不算在内。
          \item 非正式渊源。较弱效力，没有官方认同，可以考虑。\al{广西摆渡案}。
                \begin{enumerate}
                  \item 前提。禁止拒绝裁判。\zd{填补漏洞、辅助理解}。
                  \item 行业认同 \yjt 法理、判例；公众\blh \yjt 习惯 \bdy 习惯法、政策。
                \end{enumerate}
          \item 效力原则。
                \begin{enumerate}
                  \item \zd{同一主体}，新法优于旧法，特别法优于一般法。
                  \item 法规、规章 \yjt 国务院 \yjt 全人常。
                \end{enumerate}
        \end{enumerate}
  \item 法律部门。部门法。调整\zd{同一类社会关系}的法律规范的总称。\zd{同类规范}，对象重合看调整方法。
        \begin{enumerate}
          \item 公法、私法。涉及公共/私人权力、利益。
          \item 社会法。公私兼有。\al{经济法}。
        \end{enumerate}
  \item 法律体系。\zd{现行法、国内法}（包括国际私法）。9 大部门，中国特色社会法律体系\zd{已经建立}。
  \item 法的效力。约束力。
        \begin{enumerate}
          \item 默认狭义。指的是规范性文件的约束力。
          \item 来源和根据。同法概念争议。
          \item 对人效力。属人、属地、保护（操作难度大）、折中（属地为基础，其他为补充）。
          \item 空间效力。域内效力。全国有效、地方有效。\zd{只看制定主体，不看名称}。实施是条件问题。
          \item 时间效力。诉讼时效问题。
                \begin{enumerate}
                  \item 公布生效。
                  \item 失效/废止。明示（立法）\yjt 明文规定。默示（司法）\yjt 司法机关拒绝适用。\al{明大诰}。
                  \item 溯及力。对生效之前未决案能否适用。从旧、从新、从新兼从旧、\zd{从旧兼从轻}（最有利）。
                \end{enumerate}
        \end{enumerate}
  \item 法律关系。法锁。法律规范在\zd{调整社会关系}的过程中所形成的\zd{人与人}之间的\zd{权利和义务关系}。
        \begin{enumerate}
          \item \zd{社会关系 \yjt 法律规范 \yjt 法律关系}。
          \item 合法性。\zd{关系的建立}要符合法律规范的要求。不要求导致它的原因、行为合法。\al{打人赔偿}。
          \item 所有的法律关系都体现国家意志性。
          \item 权利、义务现实性。具体的，而非纸面上的。\zd{法律关系是法律规范的实现形式}。
          \item 分类。
                \begin{enumerate}
                  \item 调整性、保护性。合法要调整，违法要保护。\al{卖手机；违约}。
                  \item 横向、纵向。\zd{主体 + 地位是否平等}。
                  \item 单向、双向、多向（3 方主体）。权利义务是否对等。\al{赠予；买卖}。
                  \item 第 1 性、第 2 性。独立。主、从。\zd{调整（性）1 保护 2，实体 1 程序 2}。\al{担保合同的从属性}。
                \end{enumerate}
          \item 主体。同民法。
          \item 内容。权利、义务。
          \item 客体。权利、义务所指向的对象/载体。
                \begin{enumerate}
                  \item 物。合法、可控、可认识、经济价值、独立性。
                  \item 人身。活人。
                  \item 精神产品。
                  \item 行为结果。物化（行为 + 结果都不能出问题）\al{整容}；非物化，\al{按摩}。
                \end{enumerate}
          \item 产生、变更、消灭。法律事实（行为、事件）+ 法律规范 = 法律关系。
                \begin{enumerate}
                  \item 单一事实。一个事实影响一个 \hz 多个法律关系。
                  \item \zd{事实构成}。多个法律事实共同影响一个法律关系。\al{遗嘱}。
                \end{enumerate}
        \end{enumerate}
  \item 法律责任。承担不利后果。违法、违约、法律的特别规定。
        \begin{enumerate}
          \item 特点。责任法定、国家强制性。
          \item 责任竞合。一个人一个行为多个责任\zd{互相排斥}。就 3 个：\zd{侵权和违约、想象竞合、定金违约金}。
          \item 责任聚合。\al{刑事附带民事}。
          \item 归责。
          \item 免责。vs 无责。\al{爱会消失么}。不可抗力、正当防卫、紧急避险 \yjt 多选里视为免责，单选无责。
          \item vs 法律制裁。依照法律责任进行的一种惩罚。
        \end{enumerate}
\end{enumerate}
